%!TEX root = ../../Calc1Video.tex

%%~~~~~~~~~~~~~~~~~~~~~~~~~~~~~~~~~~~~~~~~~~~~~~~~~~~~~~~~~
%\begin{frame}
%\frametitle{{\color[rgb]{0.0,0.0,0.0} Introduction}}
%\BCOLS
%\BCOL{14.2cm}
%
%Why do we care about limits?
%
%\BCOLS
%\BCOL{8.0cm}
%\includegraphics[scale = 0.2]{Figures/VL_Calc1_Fig1_8.png}
%\includegraphics[scale = 0.2]{Figures/VL_Calc1_Fig1_9.png}
%\ECOL
%\BCOL{5.9cm}
%
%\ECOL
%\ECOLS
%
%\ECOL
%\BCOL{0.0cm}\ECOL
%\ECOLS
%\end{frame}

%%~~~~~~~~~~~~~~~~~~~~~~~~~~~~~~~~~~~~~~~~~~~~~~~~~~~~~~~~~
%\begin{frame}
%\frametitle{Why do we need Limits?}
%\BCOLS\BCOL{0cm}\ECOL\BCOL{15.5cm}\hrule\BCOLS\BCOL{0cm}\ECOL\BCOL{15.5cm}
%
%\vspace{0.2cm}
%
%\begin{center}
%{\setlength{\extrarowheight}{-0.0cm}
%\begin{tabular}{|ccccccccc|}
%	\hline
%	%\cline{1-1}\cline{3-3}\cline{5-5}\cline{7-7}\cline{9-9}
%	& & 
%	& &
%	\multirow{2}{*}{Differential} & &
%	& &
%	Solving \\
%	Limits & $\to$ & Derivatives & $\to$ & \multirow{2}{*}{Equations} & $\to$ & Modeling & $\to$ & Real World \\
%	&&&&&&&& Problems \\
%	\hline
%	%\cline{1-1}\cline{3-3}\cline{5-5}\cline{7-7}\cline{9-9}
%\end{tabular}}
%
%\vspace{1.0cm}
%
%%--------TABLE---------
%\setlength{\extrarowheight}{-0.25cm}
%\begin{tabular}{cp{7cm}}
%	\multirow{4}{*}{\includegraphics[scale = 0.7]{Figures/01_01_Fig_10.jpg}} & \\
%	\\
%	& Without the idea of limits, we would never be able comprehend how exact speed is computed. \\
%	\\
%	\\
%	& In the context of this discussion current speed will be referred to as instantaneous velocity.
%\end{tabular}
%%--------------------
%\end{center}
%
%\ECOL\BCOL{0cm}\ECOL\ECOLS\vspace{8.2cm}\ECOL\BCOL{0cm}\ECOL\ECOLS
%\end{frame}
%

%~~~~~~~~~~~~~~~~~~~~~~~~~~~~~~~~~~~~~~~~~~~~~~~~~~~~~~~~~
\begin{frame}
\frametitle{Why do we need Limits?}
\BCOLS\BCOL{0cm}\ECOL\BCOL{15.5cm}\hrule\BCOLS\BCOL{0cm}\ECOL\BCOL{15.5cm}
\begin{center}
The concept of a $limit$ defined Science and Engineering as we know it today!

\vspace{0.2cm}

{\setlength{\extrarowheight}{-0.0cm}
\begin{tabular}{|ccccccccc|}
	\hline
	%\cline{1-1}\cline{3-3}\cline{5-5}\cline{7-7}\cline{9-9}
	& & 
	& &
	\multirow{2}{*}{Differential} & &
	& &
	Solving \\
	Limits & $\to$ & Derivatives & $\to$ & \multirow{2}{*}{Equations} & $\to$ & Modeling & $\to$ & Real World \\
	&&&&&&&& Problems \\
	\hline
	%\cline{1-1}\cline{3-3}\cline{5-5}\cline{7-7}\cline{9-9}
\end{tabular}}

\vspace{0.5cm}

%--------TABLE---------
\setlength{\extrarowheight}{-0.25cm}
\begin{tabular}{cp{7cm}}
	\multirow{4}{*}{\includegraphics[scale = 0.7]{Figures/01_01_Fig_10.jpg}} & \\
	\\ \\
	& Without the idea of limits, we would never be able comprehend how speed is computed. \\
	\\ \\
	& Although we do not mention limits in this section at all, we will try to motivative the fact above. 
\end{tabular}
%--------------------
\end{center}

\ECOL\BCOL{0cm}\ECOL\ECOLS\vspace{8.2cm}\ECOL\BCOL{0cm}\ECOL\ECOLS
\end{frame}

%~~~~~~~~~~~~~~~~~~~~~~~~~~~~~~~~~~~~~~~~~~~~~~~~~~~~~~~~~
\begin{frame}
\frametitle{Average Velocity}
\BCOLS\BCOL{0cm}\ECOL\BCOL{15.5cm}\hrule\BCOLS\BCOL{0cm}\ECOL\BCOL{15.5cm}

\vspace{0.2cm}

\begin{center}
	Average Velocity = $\D{\frac{\text{Distance Traveled}}{\text{Time to Travel that Distance}}}$ = $\D{\frac{\text{Change in Position}}{\text{Change in Time}}}$ = $\D{\frac{\Delta\text{Position}}{\Delta\text{Time}}}$ 
\end{center}

\vspace{0.3cm}

{\setlength{\extrarowheight}{-0.2cm}
\begin{tabular}{p{0.25cm}p{13.6cm}}
	\EX & Suppose you travel 30 miles in a half an hour.\\\\
	& \multicolumn{1}{c}{Average Velocity = $\D{\frac{30 \text{ mi}}{0.5 \text{ h}} = 60}$ mph}\\
\end{tabular}}

\vspace{0.5cm}

It is important to understand that this does not necessarily imply that you drove 60mph the whole way.  In reality, there were certain points where you drove above 60mph and certain points where you drove below 60mph.  This is just your average speed.

\vspace{0.4cm}

{\setlength{\extrarowheight}{0.1cm}
\begin{tabular}{p{1.35cm}p{12.6cm}}
	Question: & Did you drive exactly 60mph at at least one time during your trip?\\
	Answer: & Intuition should tell you YES since you cannot drive both above and below 60mph without crossing through 60mph at some point. \\
\end{tabular}}

\ECOL\BCOL{0cm}\ECOL\ECOLS\vspace{8.2cm}\ECOL\BCOL{0cm}\ECOL\ECOLS
\end{frame}

%~~~~~~~~~~~~~~~~~~~~~~~~~~~~~~~~~~~~~~~~~~~~~~~~~~~~~~~~~
\begin{frame}
\BCOLS\BCOL{0cm}\ECOL\BCOL{15.5cm}\BCOLS\BCOL{0cm}\ECOL\BCOL{15.5cm}

\vspace{0.2cm}

Suppose we know the position of an object, $\mathcal{O}$, as a function of time.  Lets call it $s(t)$. Therefore, $s(t)$ represents the position (in meters) of $\mathcal{O}$ after some time $t$ (in seconds). 

%--------TABLE---------
\begin{center}
\setlength{\extrarowheight}{0.1cm}
\begin{tabular}{| c | c |}
	\hline
	Mathematical Expression & English Translation \\
	\hline
	$s(1)$ & 
	\ON<0>{\B{The position of $\mathcal{O}$ after 1s}}\\
	\hline
	$s(0)$ & 
	\ON<0>{\B{The initial position of $\mathcal{O}$}}\\
	\hline
	\ON<0>{\B{$s(6)$}} & 
	The position of $\mathcal{O}$ after 6s\\
	\hline
	$s(7)-s(4)$ & 
	\ON<0>{\B{$\mathcal{O}$'s change position (distance traveled) between 4s and 7s}}\\
	\hline
	\ON<0>{\B{$s(3)-s(0)$}} & 
	$\mathcal{O}$'s change position (distance traveled) in the first 3s\\
	\hline
\end{tabular}
%--------------------
\end{center}

If position, $s(t)$, of an object is known then the average velocity ,$\bar{v}$, of between time $t=a$ and $t=b$ is given by:

\vspace{0.2cm}

\centering
$\D{\bar{v} = 
\ON<0>{\B{\frac{
\only<0>{\OS{\B{ change \, in \, position }}}{
\only<0>{\OBT}{ s(b)-s(a) }}}{ 
\only<0>{\US{\B{ change \, in \, time }}}{
\only<0>{\UBT}{\, b-a \,}}}}}}$	

\ECOL\BCOL{0cm}\ECOL\ECOLS\vspace{8.2cm}\ECOL\BCOL{0cm}\ECOL\ECOLS
\end{frame}

%~~~~~~~~~~~~~~~~~~~~~~~~~~~~~~~~~~~~~~~~~~~~~~~~~~~~~~~~~
\begin{frame}
\BCOLS\BCOL{0cm}\ECOL\BCOL{15.5cm}\BCOLS\BCOL{0cm}\ECOL\BCOL{15.5cm}

\vspace{0.2cm}

{\setlength{\extrarowheight}{-0.28cm}
\begin{tabular}{p{0.25cm}p{13.6cm}}
	\EX & Consider a rock that is launched vertically upward. Let $\D{s(t) = -16t^2+96t}$ be the position (feet) of the rock after t seconds. \\
	& \\
	(a) & What is the position of the rock after 1 s? \\
	\\
	& \ON<0>{\B{$\D{s(1) = -16(1)^2+96(1) = -16+96 = 80 \, ft}$}} \\
	& \\
	(b) & What was average velocity of the rock between 1 s and 3 s? \\
	\\
	& \ON<0>{\B{$\D{\bar{v}_{1 \to 3} = \frac{s(3)-s(1)}{3-1} = \frac{[-16(3)^2+96(3)]-[-16(1)^2+96(1)]}{2} = \frac{144-80}{2} = \frac{64}{2} = 32 \, ft/s}$}} \\
	& \\
	(c) & What was average velocity of the rock between 1 s and 2 s? \\
	\\
	& \ON<0>{\B{$\D{\bar{v}_{1 \to 2} = \frac{s(2)-s(1)}{2-1} = \frac{[-16(2)^2+96(2)]-[-16(1)^2+96(1)]}{1} = \frac{128-80}{1} = 48 \, ft/s}$}} \\
	& \\
	(d) & What was average velocity of the rock between 1 s and $k$ s? (assume $k>1$)\\
	\\
	& \ON<0>{\B{$\D{\bar{v}_{1 \to k} = \frac{s(k)-s(1)}{k-1} = \frac{-16k^2+96k - 80}{k-1} \, ft/s}$}} \\
\end{tabular}}

\vspace{0.4cm}

Now, lets see what this looks like graphically...

\ECOL\BCOL{0cm}\ECOL\ECOLS\vspace{8.2cm}\ECOL\BCOL{0cm}\ECOL\ECOLS
\end{frame}

%~~~~~~~~~~~~~~~~~~~~~~~~~~~~~~~~~~~~~~~~~~~~~~~~~~~~~~~~~
\begin{frame}
\BCOLS\BCOL{0cm}\ECOL\BCOL{15.5cm}\BCOLS\BCOL{0cm}\ECOL\BCOL{15.5cm}
\centering

Position (height) of rock @ time $t = 1, 2, \& \, 3$ 
	
\vspace{0.2cm}
	
\includegraphics[scale = 0.36]{Figures/01_01_Fig_03.png}

\ECOL\BCOL{0cm}\ECOL\ECOLS\vspace{8.2cm}\ECOL\BCOL{0cm}\ECOL\ECOLS
\end{frame}

%~~~~~~~~~~~~~~~~~~~~~~~~~~~~~~~~~~~~~~~~~~~~~~~~~~~~~~~~~
\begin{frame}
\BCOLS\BCOL{0cm}\ECOL\BCOL{15.5cm}\BCOLS\BCOL{0cm}\ECOL\BCOL{15.5cm}

\centering
Average Velocity for 2 different time intervals $[1,3] \, \& \, [1,2]$
	
\vspace{0.2cm}

\begin{columns}
\begin{column}{7.5cm}
	\centering
	\includegraphics[scale = 0.27]{Figures/01_01_Fig_01.png}

	time interval: [1,3]
\end{column}
\begin{column}{7.0cm}
	\centering
	\includegraphics[scale = 0.27]{Figures/01_01_Fig_02.png}
	
	time interval: [1,2]
\end{column}
\end{columns}

\ECOL\BCOL{0cm}\ECOL\ECOLS\vspace{8.2cm}\ECOL\BCOL{0cm}\ECOL\ECOLS
\end{frame}

%~~~~~~~~~~~~~~~~~~~~~~~~~~~~~~~~~~~~~~~~~~~~~~~~~~~~~~~~~
\begin{frame}
\frametitle{Secant Lines}
\BCOLS\BCOL{0cm}\ECOL\BCOL{15.5cm}\hrule\BCOLS\BCOL{0cm}\ECOL\BCOL{15.5cm}

\begin{center}
{\setlength{\extrarowheight}{-0.4cm}
\begin{tabular}{|p{0.6cm}p{13.6cm}|}
	\hline
	\hline
	&\\
	\DEF & Let $(x_1,y_1)$ and $(x_2,y_2)$ be any two points on a curve.  The line passing through $(x_1,y_1)$ and $(x_2,y_2)$ is called a {\bf secant line.} \\
	&\\
	\hline
	\hline
\end{tabular}}
\end{center}

\centering
%--------TABLE---------
\setlength{\extrarowheight}{0.2cm}
\begin{tabular}{cp{0.2cm}c}
	\multirow{2}{*}{
	\only<1>{\includegraphics[scale = 0.27]{Figures/01_01_Fig_04.png}}
	\only<0>{\includegraphics[scale = 0.27]{Figures/01_01_Fig_05.png}}} & & 
	The slope of the secant line is \\
	&&\\
	&& 
	$\D{m_S = \frac{\Delta y}{\Delta x} = \frac{y_2-y_1}{x_2-x_1}}$\\
\end{tabular}
%--------------------	

\vspace{0.3cm}

\ECOL\BCOL{0cm}\ECOL\ECOLS\vspace{8.2cm}\ECOL\BCOL{0cm}\ECOL\ECOLS
\end{frame}

%~~~~~~~~~~~~~~~~~~~~~~~~~~~~~~~~~~~~~~~~~~~~~~~~~~~~~~~~~
\begin{frame}
\BCOLS\BCOL{0cm}\ECOL\BCOL{15.5cm}\BCOLS\BCOL{0cm}\ECOL\BCOL{15.5cm}

\centering

{\setlength{\extrarowheight}{-0.3cm}
\begin{tabular}{|p{1.2cm}p{11.7cm}|}
	\hline
	\hline
	&\\
	{\bf Concept}: & The slope of any secant line of a position curve is an average velocity\\
	&\\
	\hline
	\hline
\end{tabular}}

\vspace{0.2cm}

%--------TABLE---------
\begin{tabular}{p{5.2cm}p{1.3cm}c}
	\includegraphics[scale = 0.27]{Figures/01_01_Fig_28.png} & & 
	%\includegraphics[scale = 0.27]{Figures/01_01_Fig_06.png} 
	\\
	\multicolumn{1}{c}{position vs time} & & \\
\end{tabular}
%--------------------	

\vspace{0.3cm}

\ECOL\BCOL{0cm}\ECOL\ECOLS\vspace{8.2cm}\ECOL\BCOL{0cm}\ECOL\ECOLS
\end{frame}

%~~~~~~~~~~~~~~~~~~~~~~~~~~~~~~~~~~~~~~~~~~~~~~~~~~~~~~~~~
\begin{frame}
\BCOLS\BCOL{0cm}\ECOL\BCOL{15.5cm}\BCOLS\BCOL{0cm}\ECOL\BCOL{15.5cm}

{\setlength{\extrarowheight}{-0.3cm}
\begin{tabular}{p{0.25cm}p{14.6cm}}
	\EX & Suppose $(2,-2)$ and $(6,1)$ are two points on the graph of some function $f(x)$.  \\
	&\\
	(a) & Find the slope ($m_S$) of the secant line ($S$) between these two points \\
	&\\
	& 
	\ON<0>{\B{$\D m_S = \frac{1-(-2)}{2-6} = -\frac{3}{4}$}}\\
	&\\
	(b) & Under what condition could this slope above be considered an average velocity? \\
	&\\
	& 
	\ON<0>{\B{If $f(x)$ was a position function of $x$ (time)}}\\
\end{tabular}}

\vspace{0.4cm}

{\setlength{\extrarowheight}{-0.08cm}
\begin{tabular}{p{0.25cm}p{14.6cm}}
	\EX & Suppose the initial position of object was 5ft and its average velocity was 4ft/s during the first 8s.  Find the equation of the secant line of $s(t)$ over the time interval $[0,8]$ (secs).\\
	&\\
	& To find the equation of any line, we need the
	\ON<0>{\B{ slope, $m$, and a point, $(x_1,y_1)$, on the line}} which we plug into: 
	\ON<0>{\B{$y-y_1=m(x-x_1)$.  Turns out that we are given both: }} \\
	&\\
	& Since the average velocity = 
	\ON<0>{\B{ slope of the secant line, then $m = 4$.}} \\
	&\\
	& Since the initial position is 5ft, then 
	\ON<0>{\B{$s(0)=5$ or $(x_1,y_1) = (0,5)$. }} \\
	&\\
	& Therefore, the equation of the secant line is: 
	\ON<0>{\B{$y-5=4(x-0)$ or $y=4x+5$}} \\
\end{tabular}}

\ECOL\BCOL{0cm}\ECOL\ECOLS\vspace{8.2cm}\ECOL\BCOL{0cm}\ECOL\ECOLS
\end{frame}

%~~~~~~~~~~~~~~~~~~~~~~~~~~~~~~~~~~~~~~~~~~~~~~~~~~~~~~~~~
\begin{frame}
\frametitle{Instantaneous Velocity (or just Velocity)}
\BCOLS\BCOL{0cm}\ECOL\BCOL{15.5cm}\hrule\BCOLS\BCOL{0cm}\ECOL\BCOL{15.5cm}

\vspace{0.1cm}

We now know that the average velocity can be interpreted as the slope of some secant line.  But what about the exact velocity at a point?  It turns out that the velocity is found from the slope of the tangent line. 

\vspace{0.6cm}

\begin{columns}
\begin{column}{6.5cm}
	\includegraphics[scale = 0.2]{Figures/01_01_Fig_26.png}
\end{column}
\begin{column}{7.5cm}
	{\setlength{\extrarowheight}{-0.1cm}
\begin{tabular}{p{0.65cm}p{5.6cm}}
	\UL{Note}: & It is important to understand that it is not tangent or secant lines themselves that we are interested in, but rather the slope of these lines. \\
\end{tabular}}
\end{column}
\end{columns}
	 
\ECOL\BCOL{0cm}\ECOL\ECOLS\vspace{8.2cm}\ECOL\BCOL{0cm}\ECOL\ECOLS
\end{frame}

%~~~~~~~~~~~~~~~~~~~~~~~~~~~~~~~~~~~~~~~~~~~~~~~~~~~~~~~~~
\begin{frame}
\frametitle{Tangent Lines}
\BCOLS\BCOL{0cm}\ECOL\BCOL{15.5cm}\hrule\BCOLS\BCOL{0cm}\ECOL\BCOL{15.5cm}

To get an idea of what a tangent line is, imagine zooming into one spot of a graph as shown below.  Notice that in each box the graph appears to be a line.  If you extend the line you see in the box near $P$, you will obtain the tangent line of the function at $P$.
  
\vspace{0.2cm}

\begin{columns}
\begin{column}{8cm}
	 \includegraphics[scale = 0.22]{Figures/01_01_Fig_24.png}
\end{column}
\begin{column}{7.6cm}
{\it Important Observations:}

\vspace{0.4cm}

{\setlength{\extrarowheight}{-0.25cm}
\begin{tabular}{p{0.05cm}p{6.5cm}}
	$\triangleright$ & At each $x$ value, the graph has a $\hspace{1.0cm}$ different tangent line. \\
	&\\
	$\triangleright$ & Each tangent line has a slope, which measures how steep the line is. \\
	&\\
	$\triangleright$ & At any point on the graph (e.g. $P$), the function and the tangent line share the same steepness. \\
	&\\
	&Therefore...\\
\end{tabular}}	
\end{column}
\end{columns}
\!\!
\begin{center}
{\setlength{\extrarowheight}{-0.5cm}
\begin{tabular}{|p{1.2cm}p{13.1cm}|}
	\hline
	\hline
	&\\
	{\bf Concept}: & The slope of the tangent line at $P$ tells us the steepness of the function at $P$.\\
	&\\
	\hline
	\hline
\end{tabular}}
\end{center}

\ECOL\BCOL{0cm}\ECOL\ECOLS\vspace{8.2cm}\ECOL\BCOL{0cm}\ECOL\ECOLS
\end{frame}

%~~~~~~~~~~~~~~~~~~~~~~~~~~~~~~~~~~~~~~~~~~~~~~~~~~~~~~~~~
\begin{frame}
\BCOLS\BCOL{0cm}\ECOL\BCOL{15.5cm}\BCOLS\BCOL{0cm}\ECOL\BCOL{15.5cm}

In general, we need 2 points on a line to find the slope of the line.  However, tangent lines only cross through 1 point, $P$.  It turns out that the slope of the tangent line can be approximated from the the slope of a secant line that is very close to the tangent line.
  
\vspace{0.2cm}

Consider adding a point $Q_1$ to the curve and draw a secant line through $P$ and $Q_1$.  We can compute the slope of this secant line, but clearly, this won't give us the slope of the tangent line.  Now, imagine that we added more points $Q_2$, $Q_3$, ... closer to $P$.  What will happen to the blue secant lines as the points keep getting closer to $P$?

\vspace{0.1cm}

\begin{columns}
\begin{column}{7.2cm}
\vspace{0.2cm}
	\includegraphics[scale = 0.19]{Figures/01_01_Fig_13.png}
\end{column}
\begin{column}{7.5cm}
	\includegraphics[scale = 0.19]{Figures/01_01_Fig_22.png}
\end{column}
\end{columns}

\vspace{0.1cm}

It should be clear that the secant lines will eventually be very close to the tangent line.  Thus, the slopes of secant lines will also eventually be very close to the slope of the tangent line.
 
%\begin{center}
%	\only<1>{\includegraphics[scale = 0.22]{Figures/01_01_Fig_12.png}}
%	\only<1>{\includegraphics[scale = 0.22]{Figures/01_01_Fig_14.png}}
%	\only<0>{\includegraphics[scale = 0.26]{Figures/01_01_Fig_13.png}}
%	\only<0>{\includegraphics[scale = 0.26]{Figures/01_01_Fig_19.png}}
%	\only<0>{\includegraphics[scale = 0.26]{Figures/01_01_Fig_20.png}}
%	\only<0>{\includegraphics[scale = 0.26]{Figures/01_01_Fig_22.png}}
%\end{center}

\ECOL\BCOL{0cm}\ECOL\ECOLS\vspace{8.2cm}\ECOL\BCOL{0cm}\ECOL\ECOLS
\end{frame}

%~~~~~~~~~~~~~~~~~~~~~~~~~~~~~~~~~~~~~~~~~~~~~~~~~~~~~~~~~
\begin{frame}
\BCOLS\BCOL{0cm}\ECOL\BCOL{15.5cm}\BCOLS\BCOL{0cm}\ECOL\BCOL{15.5cm}

{\setlength{\extrarowheight}{-0.3cm}
\begin{tabular}{p{0.25cm}p{14.6cm}}
	\EX & Given the position function, $\D s(t)=5e^{t}$ (meters), estimate the instantaneous velocity at $t=2.5$s to within 2 decimal places.  \\
	&\\
	& In order to estimate the velocity(slope of the tangent line) at $t=2.5$s we need to use a sequence of average velocities (slopes of secant lines) near $t=2.5$s (just like above). Lets compute the average velocities on the following intervals: \\
	& 
	%--------TABLE---------
	\multicolumn{1}{c}{
{\setlength{\extrarowheight}{-0.4cm}
\begin{tabular}{| c | l |}
	\hline
	&\\
	$[2.5,2.51]$ & $\D \frac{s(2.51)-s(2.5)}{2.51-2.5} = \frac{5e^{2(2.51)}-5e^{2(2.5)}}{0.01} = 61.2181$\\
	&\\
	\hline
	&\\
	$[2.5,2.51]$ & $\D \frac{s(2.501)-s(2.5)}{2.501-2.5} = \frac{5e^{2(2.501)}-5e^{2(2.5)}}{0.001} = 60.9429$\\
	&\\
	\hline
	&\\
	$[2.5,2.5001]$ & $\D \frac{s(2.5001)-s(2.5)}{2.5001-2.5} = \frac{5e^{2(2.5001)}-5e^{2(2.5)}}{0.0001} = \UL{60.91}55$\\
	&\\
	\hline
	&\\
	$[2.5,2.50001]$ & $\D \frac{s(2.50001)-s(2.5)}{2.50001-2.5} = \frac{5e^{2(2.50001)}-5e^{2(2.5)}}{0.00001} = \UL{60.91}28$\\
	&\\
	\hline
\end{tabular}}}\\
	&\\
	& Notice that the average velocities appear to be converging some fixed value as the intervals continue to shrink. The value that we would eventually reach is the instantaneous velocity at $t=2.5$s.  For this example, we just estimate the this value as \fbox{$60.91$m/s} \\
\end{tabular}}

\ECOL\BCOL{0cm}\ECOL\ECOLS\vspace{8.2cm}\ECOL\BCOL{0cm}\ECOL\ECOLS
\end{frame}


%
%%~~~~~~~~~~~~~~~~~~~~~~~~~~~~~~~~~~~~~~~~~~~~~~~~~~~~~~~~~
%\begin{frame}
%\BCOLS\BCOL{0cm}\ECOL\BCOL{15.5cm}\BCOLS\BCOL{0cm}\ECOL\BCOL{15.5cm}
%
%Below is the idea for approximating the slope of a tangent line at a point $P$.  Let $m_T$ be the slope of the tangent line which we want to approximate.
%%-------LIST---------
%\begin{itemize}
%	\item Label the $x$-coord of $P$ as $x$, so $P$ is $(x,f(x))$
%	\item Label the $x$-coord of $Q$ as $x+h$ \, ($h$ is some value that we control), so $Q$ is $(x+h,f(x+h))$
%	\item Note: We can make $Q$ move toward $P$ by decreasing the value $h$
%	\item Compute the slope the secant line between $P$($x$,$f(x)$) \& $Q$($x+h$,$f(x+h)$) and call it $m_S$ 
%\end{itemize}
%%-------------------
%
%\vspace{0.1cm}
%
%\centering
%
%\begin{columns}
%\begin{column}{0.0cm}
%\end{column}
%\begin{column}{2.0cm}
%	%--------EQN----------
%	\begin{align*}
%		m_S & = \frac{y_2-y_1}{x_2-x_1} \\
%		    & = \frac{f(x+h)-f(x)}{x+h-x} \\
%		    & = \frac{f(x+h)-f(x)}{h} \\
%	\end{align*}
%	%----------------------
%\end{column}
%\begin{column}{10.5cm}
%	\includegraphics[scale = 0.23]{Figures/01_01_Fig_27.png}
%\end{column}
%\end{columns}
%
%\vspace{0.2cm}
%
%{\it Idea}: \, If $h$ is small, then $S$ is close to $T$, and thus slope of $S$ is close to the slope of $T$
%
%In calculus, we say that as 
%$\only<0>{\OS{\B{ h \to 0  }}}{\only<0>{\OBT}{ \text{$h$ approaches $0$} }}$, 
%$\only<0>{\OS{\B{ S \to T  }}}{\only<0>{\OBT}{ \text{$S$ approaches $T$} }}$, and 
%$\only<0>{\OS{\B{ m_S \to m_T  }}}{\only<0>{\OBT}{ \text{$m_S$ approaches $m_T$} }}$ 
% 
%\ECOL\BCOL{0cm}\ECOL\ECOLS\vspace{8.2cm}\ECOL\BCOL{0cm}\ECOL\ECOLS
%\end{frame}
%
%%~~~~~~~~~~~~~~~~~~~~~~~~~~~~~~~~~~~~~~~~~~~~~~~~~~~~~~~~~
%\begin{frame}
%\BCOLS\BCOL{0cm}\ECOL\BCOL{15.5cm}\BCOLS\BCOL{0cm}\ECOL\BCOL{15.5cm}
%
%{\setlength{\extrarowheight}{-0.1cm}
%\begin{tabular}{p{0.25cm}p{13.6cm}}
%	\EX & Approximate the slope of the tangent line(red) of $f(x) = (
%	\!\!\BLINE{\node[anchor=base](n1){$x$};}\!\!-1.8)^2$ at  
%	\!\!\BLINE{\node[anchor=base](n2){$x=1$};}\!\!\\
%	&\\
%	& Lets set up the 2 points of the movable secant line(blue). \\
%	& Pt 1:  $\big(
%	\BLINE{\node[anchor=base](n3){1};},\, 
%	\only<0>{\OS{\B{ f(1) }}}{
%	\only<0>{\OBT}{ [1-1.8]^2 }} \, \big) = (1,0.64)$  \hspace{0.5cm} Pt 2:  $\big( 1+h, 
%	\only<0>{\OS{\B{ f(1+h) }}}{\only<0>{\OBT}{ [(1+h)-1.8]^2 }} \, \big) = (1+h,[h-0.8]^2)$ \\
%\end{tabular}}
%
%\tikzstyle{na} = [baseline=-2.5ex]
%\begin{tikzpicture}[overlay]
%	\ON<0>{\ARROWB (n2.\DS) \BENDL (n3.\DS);}
%\end{tikzpicture}
%
%\begin{tikzpicture}
%
%\begin{axis}[
%    %title={Temperature dependence of CuSO$_4\cdot$5H$_2$O solubility},
%%    xlabel={$x$},
%%    ylabel={$y$},
%    axis lines = center,
%    xmin=0, xmax=2.6,
%    ymin=-0.5, ymax=2,
%    xtick={1,2},
%    ytick={1,2},
%    xticklabels={1,$1+h$},%<--Here
%    %xlabel style={yshift=-1cm},
%    %legend pos=north west,
%    %ymajorgrids=true,
%    %grid style=dashed,
%]
%\node[label={180:{(1,0.64)}},circle,fill,inner sep=1pt] at (axis cs:1,0.64) {};
%\node[label={90:{$(1+h,[h-0.8]^2)$}},circle,fill,inner sep=1pt] at (axis cs:2,0.04) {};
%   	
%\addplot[domain=0:3,
%    color=black, thick
%    ] {(x-1.8)^2};
%
%\addplot[domain=0.3:1.8,
%    color=red
%    ] {-1.6*x+2.24};
%    
%\addplot[color=blue]
%	coordinates {
%    	(1,0.64)(2,0.04)
%    };
%
%\end{axis}
%\end{tikzpicture}
%%
%
%
%%\begin{tikzpicture}%[scale=0.75]
%%    \begin{axis}[
%%    title={Temperature dependence of CuSO$_4\cdot$5H$_2$O solubility},
%%    ylabel={Solubility [g per 100 g water]},
%%    xmin=0.2, xmax=1
%%%%    ymin=0, ymax=120,
%%%    xtick={0,20,40,60,80,100},
%%%    ytick={0,20,40,60,80,100,120},
%%%    legend pos=north west,
%%%    ymajorgrids=true,
%%%    grid style=dashed,
%%]
%%    	\addplot[color=blue] {1/x};
%%%		%\legend{$\sin(x)$,$\cos(x)$,$x^2$}
%%%		\addplot[mark=square,color=red]
%%%    	coordinates {
%%%    		(1,1)
%%%    	};
%%    \end{axis}
%%\end{tikzpicture}
%
%This process of using known values to approach a unknown value is the basis of limits and is one of the most fundamental ideas of mathematics and consequently all of science!
%\ECOL\BCOL{0cm}\ECOL\ECOLS\vspace{8.2cm}\ECOL\BCOL{0cm}\ECOL\ECOLS
%\end{frame}
%
%\newcounter{r}
%\newcommand{\escalar}[1]{
%\setcounter{r}{#1 * #1 * #1}
%}
%%
%\newcounter{m}
%\setcounter{m}{0}
%\newcounter{mc}
%%
%%
%%
%%~~~~~~~~~~~~~~~~~~~~~~~~~~~~~~~~~~~~~~~~~~~~~~~~~~~~~~~~~
%
%\begin{frame}[fragile]{Animated Integral}
%\begin{center}
%\begin{animateinline}[loop, poster = first, controls, palindrome]{2}
%\whiledo{\them < 21}{
%    \begin{tikzpicture}[scale=1.25]
%    \draw[red,thick,<->] (-1,1) parabola bend (0,0) (2.1,4.41)
%        node[below right] {$y=x^2$};
%    \draw[loosely dotted] (-1,0) grid (4,4);
%    %\path[use as bounding box] (-2,-1) rectangle (5,5);
%    \draw[->] (-0.2,0) -- (4.25,0) node[right] {$x$};
%    \draw[->] (0,-0.25) -- (0,4.25) node[above] {$y$};
%    \foreach \x/\xtext in {1/1, 2/2, 3/3}
%    \draw[shift={(\x,0)}] (0pt,2pt) -- (0pt,-2pt) node[below] {$\xtext$};
%    \foreach \y/\ytext in {1/1, 2/2, 3/3, 4/4}
%    \draw[shift={(0,\y)}] (2pt,0pt) -- (-2pt,0pt) node[left] {$\ytext$};
%%
%    \setcounter{mc}{\value{m}*\value{m}}
%    \shade[top color=blue,bottom color=gray!50]
%        (0,0) parabola (0.1*\them,0.01*\themc) |- (0,0);
%    \escalar{\them}
%    \draw (3cm,2pt) node[above]
%        {$\displaystyle\int\limits_0^{\them/10} \!\!x^2\mathrm{d}x = 
%            \displaystyle\frac{\ther}{3000}$};
%    \draw[fill=orange,color=orange] (0.1*\them,0.01*\themc) circle (0.5pt);
%    \end{tikzpicture}
%    %
%    \stepcounter{m}
%    \ifthenelse{\them < 21}{
%            \newframe
%    }{
%        \end{animateinline}\relax % BREAK
%    }
%} % END \whiledo...
%\end{center}
%\end{frame}

